\section{Conclusion}

We presented AMT-UNet, a hybrid CNN-transformer architecture for joint brain tumor segmentation and grade classification. The method combines: (1) an adaptive masked transformer bottleneck that focuses self-attention on tumor-relevant patches, reducing complexity while improving context modeling, (2) an encoder-decoder with instance normalization, residual connections, CBAM attention, and multi-scale fusion for robust feature extraction, (3) ROI-guided multi-task learning with gradient stopping to prevent task interference, and (4) a multi-component loss (Dice, focal, IoU, boundary) with deep supervision for effective optimization.

Experiments on BraTS 2020 demonstrate competitive performance for both segmentation and classification with 37M parameters. Key design choices include: instance normalization for scanner independence, strong focal focusing ($\gamma=3.0$) for extreme class imbalance, IoU weight (2.0) for metric alignment, and soft masking for efficient global attention.

\subsection{Limitations and Future Work}

Current limitations suggest several directions: (1) \textbf{3D extension} to leverage inter-slice context through hierarchical or sparse processing, (2) \textbf{uncertainty quantification} via Bayesian methods or ensembles for calibrated confidence, (3) \textbf{multi-center validation} with domain adaptation for improved generalization, (4) \textbf{model compression} through distillation or pruning for deployment, (5) \textbf{interpretability} beyond attention visualization for clinical trust, (6) \textbf{extended multi-task learning} for survival prediction and molecular subtyping, (7) \textbf{active learning} to reduce annotation costs, and (8) \textbf{few-shot adaptation} to new tumor types or modalities.

AMT-UNet demonstrates that principled combination of CNNs and transformers, guided by domain-specific requirements, achieves effective medical image analysis. The adaptive masking mechanism and ROI-guided multi-task design provide practical solutions to computational and optimization challenges in transformer-based medical imaging.

\begin{credits}
\subsubsection{\ackname}
\textbf{[TODO: Add acknowledgments, for example: This study was funded by X (grant number Y). We thank the BraTS 2020 challenge organizers for providing the dataset.]}

\subsubsection{\discintname}
\textbf{[TODO: Declare competing interests, for example: The authors have no competing interests to declare that are relevant to the content of this article.]}
\end{credits}
